\chapter{Linear Algebra Reference}\label{app:a}

\section*{Purpose and learning objectives}
This appendix provides the calculations used repeatedly in the book. Readers should be able to check dimensions, distinguish a subspace from a list of coordinates, compute a projection, interpret a singular value, and differentiate a quadratic loss. Chapter~\ref{ch:01} develops dimensionality reduction from these ingredients. \citet{ev_axler2024} is a more extensive reference.

\section{Vectors, maps, and coordinates}
A real vector space is a set with addition and real scalar multiplication satisfying the familiar distributive, associative, and identity rules. Its elements need not be arrows: functions and polynomials also form vector spaces. A subspace contains zero and is closed under addition and scalar multiplication.

Vectors \(v_1,\ldots,v_k\) are linearly independent if
\[
\sum_{j=1}^k a_jv_j=0\quad\Longrightarrow\quad a_1=\cdots=a_k=0.
\]
Their span is the set of all such linear combinations. A basis is an independent spanning list; its length is the dimension. Coordinates are coefficients in a chosen basis, not properties of a vector independently of that basis.

A linear map \(A:\R^p\to\R^m\) satisfies \(A(ax+by)=aAx+bAy\) and is represented by an \(m\times p\) matrix. Its columns are the images of the standard basis vectors. Its image and null space are
\[
\operatorname{im}A=\{Ax:x\in\R^p\},\qquad
\ker A=\{x:Ax=0\}.
\]
The rank is \(\dim(\operatorname{im}A)\); the nullity is \(\dim(\ker A)\).

\begin{theorem}[Rank--nullity]
For \(A:\R^p\to\R^m\), \(\rank A+\dim\ker A=p\).
\end{theorem}
\begin{proof}[Proof sketch]
Extend a basis \(v_1,\ldots,v_r\) of the kernel to a basis of the domain. The images of the added basis vectors span the image; they are independent because a relation among their images would place a nontrivial combination of the added vectors in the kernel. There are \(p-r\) such vectors.
\end{proof}

\begin{example}
For
\[
A=\begin{pmatrix}1&1&0\\0&0&1\end{pmatrix},
\]
\(Ax=(x_1+x_2,x_3)^\top\). The image is all of \(\R^2\), and the null space is \(\operatorname{span}\{(1,-1,0)^\top\}\). A two-dimensional readout cannot recover which of the first two coordinates supplied their sum.
\end{example}

\section{Inner products, norms, and orthogonality}
The Euclidean inner product is \(\ip{x}{y}=x^\top y\). Its induced norm is \(\norm{x}_2=\sqrt{x^\top x}\); vectors are orthogonal when their inner product is zero. The Cauchy--Schwarz inequality states
\[
\abs{x^\top y}\leq\norm x_2\norm y_2.
\]
For \(y\ne0\), expand \(\norm{x-ty}^2\geq0\) and minimize over \(t\) to obtain the inequality. This proof also explains the projection coefficient \(t=x^\top y/(y^\top y)\).

An orthonormal matrix \(Q\in\R^{p\times k}\) satisfies \(Q^\top Q=I_k\). Unless \(k=p\), generally \(QQ^\top\ne I_p\). Instead \(QQ^\top\) is the projector onto the column space of \(Q\):
\[
P^2=P,\qquad P^\top=P,\qquad
x=Px+(I-P)x,\qquad (Px)^\top(I-P)x=0.
\]
The last identity gives the Pythagorean decomposition \(\norm x^2=\norm{Px}^2+\norm{(I-P)x}^2\).

\section{Least squares and orthogonal projection}
Given \(A\in\R^{n\times k}\) and \(y\in\R^n\), minimize
\[
L(w)=\frac12\norm{Aw-y}_2^2.
\]
Expanding the quadratic and differentiating gives
\[
\nabla_wL=A^\top(Aw-y).
\]
At a minimizer, the residual is orthogonal to every column of \(A\). If the columns are independent,
\[
\widehat w=(A^\top A)^{-1}A^\top y,\qquad
\widehat y=A(A^\top A)^{-1}A^\top y.
\]
The fitted vector \(\widehat y\) is unique even when the coefficients are not. For a rank-deficient design, a pseudoinverse selects the minimum-Euclidean-norm solution. In numerical work, QR or SVD methods are generally preferable to forming the inverse of \(A^\top A\), whose conditioning can be substantially worse.

\begin{example}[Projection with numbers]
Project \(y=(2,0)^\top\) onto the line spanned by \(a=(1,1)^\top\). The coefficient is \(a^\top y/a^\top a=1\), so \(\widehat y=(1,1)^\top\). The residual \((1,-1)^\top\) is orthogonal to \(a\), and the squared norms satisfy \(4=2+2\).
\end{example}

\section{Eigenvalues and the spectral theorem}
For a square matrix \(A\), an eigenvector \(v\ne0\) satisfies \(Av=\lambda v\). An eigenvalue is a scalar \(\lambda\) for which such a vector exists. General real matrices can have complex eigenvalues, defective eigenspaces, and nonorthogonal eigenvectors.

\begin{theorem}[Real spectral theorem]
Every real symmetric matrix \(A=A^\top\) has an orthonormal eigenbasis. Thus
\[
A=Q\Lambda Q^\top,
\]
where \(Q^\top Q=I\) and \(\Lambda\) is real diagonal.
\end{theorem}
Symmetry is essential. For a proof sketch, maximize \(x^\top Ax\) on the compact unit sphere. A Lagrange multiplier gives an eigenvector. Its orthogonal complement is invariant under \(A\); repeat on that complement. The result implies that the extrema of the Rayleigh quotient \(x^\top Ax/(x^\top x)\) are the smallest and largest eigenvalues.

For \(A=\begin{pmatrix}2&1\\1&2\end{pmatrix}\), the normalized eigenvectors \((1,1)^\top/\sqrt2\) and \((1,-1)^\top/\sqrt2\) have eigenvalues \(3\) and \(1\). The same eigenvectors diagonalize the quadratic form.

\section{Singular value decomposition}
Every real \(n\times p\) matrix has an SVD
\[
A=U\Sigma V^\top,
\]
with orthogonal \(U,V\) and nonnegative singular values \(\sigma_1\geq\cdots\geq\sigma_r>0\) on the rectangular diagonal of \(\Sigma\). Here \(r=\rank A\). The right singular vectors \(v_j\) satisfy \(A^\top A v_j=\sigma_j^2v_j\), and \(Av_j=\sigma_ju_j\).

One construction diagonalizes the positive semidefinite matrix \(A^\top A\), then sets \(u_j=Av_j/\sigma_j\) for each positive singular value and completes orthonormal bases. Unlike eigenvalue decomposition, the SVD applies to rectangular matrices and always exists.

The Frobenius norm is
\[
\norm A_F^2=\sum_{ij}A_{ij}^2=\tr(A^\top A)=\sum_j\sigma_j^2.
\]
The operator norm induced by the Euclidean norm is
\[
\norm A_2=\max_{\norm x_2=1}\norm{Ax}_2=\sigma_1.
\]
The pseudoinverse is \(A^+=V\Sigma^+U^\top\), where positive singular values are inverted and zero singular values remain zero. Tiny singular values make inversion sensitive; a numerical threshold or regularization is then a modeling choice.

\section{Positive definiteness and quadratic forms}
A real symmetric matrix \(M\) is positive semidefinite if \(x^\top Mx\geq0\) for all \(x\), and positive definite if the inequality is strict for every nonzero \(x\). By the spectral theorem, these conditions mean respectively nonnegative and positive eigenvalues.

Covariance matrices are positive semidefinite because
\[
x^\top\Cov(Y)x=\Var(x^\top Y)\geq0.
\]
They need not be invertible. If one measured channel is an exact linear combination of others, there is a zero-variance direction. A Mahalanobis distance with \(M=\Sigma^{-1}\) is defined as
\[
d_M(x,y)^2=(x-y)^\top M(x-y).
\]
Positive definiteness makes this a genuine squared metric distance; a singular \(M\) can assign zero distance to distinct points.

\section{Matrix calculus and shape checks}
For a scalar function \(L(W)\), define its gradient through
\[
\dd L=\tr\big((\nabla_WL)^\top\dd W\big).
\]
This convention makes the gradient the same shape as \(W\). The identities
\[
\dd(AB)=(\dd A)B+A\,\dd B,\quad
\dd(W^\top)=(\dd W)^\top,\quad
\tr(AB)=\tr(BA)
\]
are sufficient for many derivations. Cyclic permutation inside a trace is allowed when products are defined; arbitrary rearrangement is not.

For \(L(W)=\frac12\norm{XW-Y}_F^2\), let \(E=XW-Y\). Then
\[
\dd L=\tr(E^\top X\,\dd W)
=\tr\big((X^\top E)^\top\dd W\big),
\qquad \nabla_WL=X^\top(XW-Y).
\]
If \(X\) is \(n\times p\) and \(W\) is \(p\times q\), this gradient is \(p\times q\), as required.

For \(f:\R^p\to\R^m\), the Jacobian has entries \(J_{ij}=\partial f_i/\partial x_j\). For scalar \(L\), the Hessian has entries \(H_{ij}=\partial^2L/\partial x_i\partial x_j\). If the second partial derivatives are continuous, \(H\) is symmetric. For \(L(x)=\frac12x^\top Mx\), the gradient is \((M+M^\top)x/2\); it simplifies to \(Mx\) only when \(M\) is symmetric.

\section*{Exercises and quick checks}
\begin{enumerate}
\item \textbf{Introductory.} Find the rank and null space of \(\begin{pmatrix}1&2\\2&4\end{pmatrix}\).
\item \textbf{Intermediate.} Verify directly that \(P=A(A^\top A)^{-1}A^\top\) is symmetric and idempotent when \(A\) has independent columns.
\item \textbf{Advanced.} Derive the ridge solution \((A^\top A+\lambda I)^{-1}A^\top y\) and express its shrinkage in the singular-vector basis.
\end{enumerate}
\section*{Summary and further reading}
The essential checks are shapes, symmetry, and limiting cases. If \(A=I\), least squares should return \(y\); if a projection retains the full space, its residual should vanish; and a covariance quadratic form should never be negative apart from numerical error. A common misconception is that every square matrix has an orthogonal eigenbasis: this is guaranteed for real symmetric matrices, whereas general matrices require different qualifications. For a systematic development, see \citet{ev_axler2024}; return to Chapter~\ref{ch:01} for the connection to PCA.
